\subsection{Typical truncation on a physical region}

\begin{definition}[Regional typical pure vector]
\label{def:entropy_regional_typical_state}
\lean{FiniteProduct.typicalPureState,FiniteProduct.split_typicalPureState}
\leanok
\uses{def:entropy_typical_pure_state,def:finite_product_regional_entropy}
Let $V$ be a finite set of sites with finite local Hilbert spaces, and let
$\psi\in\bigotimes_{v\in V}\mathcal H_v$. For a region $X\subseteq V$, let
$\rho_X=\operatorname{tr}_{X^c}|\psi\rangle\langle\psi|$.
Given a set $E$ of its eigenvalue indices, let $P_E$ be the selected spectral
projection and let $z=\operatorname{tr}(P_E\rho_X)$.
If $z>0$, the global selected vector is
\begin{align}
 \psi_E=z^{-1/2}(P_E\otimes I_{X^c})\psi,
\end{align}
expressed in the original product basis. This vector is fixed by $X$ and $E$.
\end{definition}

\begin{theorem}[Normalization and regional marginal]
\label{thm:entropy_regional_typical_state_normalization}
\lean{FiniteProduct.norm_typicalPureState,FiniteProduct.reducedPure_typicalPureState}
\leanok
\uses{def:entropy_regional_typical_state,thm:entropy_typical_pure_state}
If $z>0$, then $\|\psi_E\|=1$ and
\begin{align}
 \operatorname{tr}_{X^c}|\psi_E\rangle\langle\psi_E|
   =z^{-1}P_E\rho_XP_E.
\end{align}
\end{theorem}
\begin{proof}
\leanok
The canonical decomposition into configurations on $X$ and $X^c$ is an isometry.
The claims follow from the normalization and first marginal of the selected
bipartite vector in these coordinates.
\end{proof}

\begin{theorem}[Marginal estimates on every disjoint physical region]
\label{thm:entropy_regional_typical_state_bound}
\lean{FiniteProduct.reducedPure_typicalPureState_le,FiniteProduct.entropy_typicalPureState_le}
\leanok
\uses{def:entropy_regional_typical_state,thm:entropy_typical_complementary_subsystem}
Let $z>0$. For every region $B$ disjoint from $X$, the same selected global
vector satisfies
\begin{align}
 \rho_{\psi_E,B}\le z^{-1}\rho_{\psi,B}.
\end{align}
If also $\|\psi\|=1$, then
\begin{align}
 S_B(\psi_E)\le \frac{S_B(\psi)}{z}.
\end{align}
These are the disjoint-region estimates of
\cite[Section~8]{OpenAI2026AreaLaw}.
\end{theorem}
\begin{proof}
\leanok
Decompose the configurations on $X^c$ into those on $B$ and those on
$(X\cup B)^c$. Reindexing the complementary basis preserves the marginal on $X$
and commutes with the selected projection. Thus the same global vector is the
three-factor selected vector in these coordinates. Taking the partial traces
over $X$ and $(X\cup B)^c$ gives precisely the actual regional marginal on $B$.
Apply the three-factor positive-semidefinite and entropy inequalities.
\end{proof}
